% Cycle versus cut: the same pairwise disagreement is detectable on C6, where
% it closes a visible cycle, and invisible on P6, where no edge closes one.
% Radii from skills/harbor-results/scripts/sheaf_consistency_radius.py
% (seed 20260816): C6 r = 3 sqrt(1 - 5/6) = 1.2247, P6 r = 0. [verified]
%
% Reader question: why does the same disagreement get caught on one topology
%   and not another?
% Claim: detection is exactly the visible-cycle condition: the disagreeing
%   overlap closes a cycle on C6 (r>0) but no edge closes one on P6 (r=0).
% Evidence: the two computed radii from sheaf_consistency_radius.py, seed
%   20260816, on six harbors numbered 1-6.
% Must distinguish: a closed cycle through the disagreement (detectable)
%   from the same disagreement on a path (invisible, r=0 exactly).
% Grammar chosen: before/after pair, two instances of the same six-node
%   graph on one grid (tpl-before-after idiom, applied to two topologies
%   rather than two times). Grammar rejected: a single graph with both edge
%   sets overlaid would hide which edge is the one difference (the missing
%   closing edge on the path).
% Five-second test: which panel's disagreeing edge closes a cycle, and what
%   is that panel's radius?
\pdfigurehue{pdgold}%
\begin{figure}[H]
\centering
\begin{tikzpicture}[pd figure,x=1cm,y=1cm]
  \def\cax{1.50}   \def\pbx{7.10}   \def\rad{1.25}
  \node[pd title,anchor=north,text width=4.0cm,align=center] at (\cax,2.85)
    {6-cycle\\closes a visible loop};
  \foreach \i in {1,...,6}
    {\node[pd state,circle,minimum size=6.2mm,inner sep=0pt] (c\i)
       at ({\cax+\rad*cos(90-60*(\i-1))},{\rad*sin(90-60*(\i-1))}) {\i};}
  \foreach \i/\j in {1/2,2/3,3/4,4/5,5/6}
    {\draw[pd rule] (c\i) -- (c\j);}
  \draw[pd breach rule] (c6) -- (c1);
  \node[pd label,anchor=north,text width=3.6cm,align=center] at (\cax,-1.75)
    {$r=1.2247$};

  \node[pd title,anchor=north,text width=4.6cm,align=center] at (\pbx,2.85)
    {6-path\\has no visible loop};
  \foreach \i in {1,...,6}
    {\node[pd state,circle,minimum size=6.2mm,inner sep=0pt] (p\i)
       at ({\pbx-2.15+(\i-1)*0.86},0.30) {\i};}
  \foreach \i/\j in {1/2,2/3,4/5,5/6}
    {\draw[pd rule] (p\i) -- (p\j);}
  \draw[pd breach rule] (p3) -- (p4);
  \node[pd label,anchor=north,text width=3.6cm,align=center] at (\pbx,-1.75)
    {$r=0$};
  % One shared legend in its own lane; labels never enter the other panel.
  \draw[pd breach rule] (2.15,-2.75)--(3.15,-2.75);
  \node[pd label,anchor=west] at (3.4,-2.75) {edge perturbation of size $3$};
\end{tikzpicture}
\caption{One-edge perturbation of size 3: the cycle residual is $1.2247$; the path residual is zero. This is an edge-cochain diagnostic. Authenticated endpoint claims can be compared directly on either topology. Analytic construction.}
\label{fig:fh-cycle-vs-cut}
\end{figure}
